\documentclass[10pt,twocolumn]{article}

\usepackage[margin=0.72in]{geometry}
\usepackage{amsmath}
\usepackage{booktabs}
\usepackage{graphicx}
\usepackage[hidelinks]{hyperref}
\usepackage{xcolor}
\usepackage{enumitem}
\setlist{nosep,leftmargin=*}
\setlength{\columnsep}{0.24in}
\setlength{\parindent}{1em}
\setlength{\parskip}{0pt}

% This is the only numeric/text data file normally edited when finalizing the report.
% Keep every placeholder until the corresponding frozen evidence cell has been inspected.
% Escape LaTeX special characters in free-text values.

\newcommand{\FGReleaseVersion}{0.1.0-rc1}
\newcommand{\FGAuthor}{Perilla-desuwa}
\newcommand{\FGReleaseRevision}{v0.1.0-rc1}
\newcommand{\FGRamseySourceRevision}{4a9a98954f95}
\newcommand{\FGJavaScriptSourceRevision}{834de4296861}
\newcommand{\FGCPAblationSourceRevision}{3c0a87122dc8}
\newcommand{\FGEvidenceBundleChecksum}{the accompanying \texttt{.sha256} release asset}
\newcommand{\FGExperimentDate}{2026-08-18}
\newcommand{\FGHostSummary}{AMD Ryzen 7 7745HX; 16 logical processors; 63.69 GiB RAM; Windows 11; Python 3.12.13}
\newcommand{\FGPlannerModel}{gpt-5.6-luna}
\newcommand{\FGReasoningEffort}{low for Ramsey and TinyShop; medium planning and low worker reasoning for JavaScript}

\newcommand{\FGRamseyRuns}{eight completed cells: slots 1/2/4/8, two repetitions each}
\newcommand{\FGRamseyOneMedian}{503.29}
\newcommand{\FGRamseyTwoMedian}{236.06}
\newcommand{\FGRamseyFourMedian}{153.48}
\newcommand{\FGRamseyEightMedian}{153.82}

\newcommand{\FGRamseyOneRepOneTime}{446.45}
\newcommand{\FGRamseyOneRepTwoTime}{560.13}
\newcommand{\FGRamseyTwoRepOneTime}{263.05}
\newcommand{\FGRamseyTwoRepTwoTime}{209.06}
\newcommand{\FGRamseyFourRepOneTime}{139.69}
\newcommand{\FGRamseyFourRepTwoTime}{167.28}
\newcommand{\FGRamseyEightRepOneTime}{170.08}
\newcommand{\FGRamseyEightRepTwoTime}{137.56}

\newcommand{\FGRamseyOneRepOneSpeedup}{1.00}
\newcommand{\FGRamseyOneRepTwoSpeedup}{0.80}
\newcommand{\FGRamseyTwoRepOneSpeedup}{1.70}
\newcommand{\FGRamseyTwoRepTwoSpeedup}{2.14}
\newcommand{\FGRamseyFourRepOneSpeedup}{3.20}
\newcommand{\FGRamseyFourRepTwoSpeedup}{2.67}
\newcommand{\FGRamseyEightRepOneSpeedup}{2.62}
\newcommand{\FGRamseyEightRepTwoSpeedup}{3.25}

\newcommand{\FGRamseyOneRepOneEfficiency}{100.0\%}
\newcommand{\FGRamseyOneRepTwoEfficiency}{79.7\%}
\newcommand{\FGRamseyTwoRepOneEfficiency}{84.9\%}
\newcommand{\FGRamseyTwoRepTwoEfficiency}{106.8\%}
\newcommand{\FGRamseyFourRepOneEfficiency}{79.9\%}
\newcommand{\FGRamseyFourRepTwoEfficiency}{66.7\%}
\newcommand{\FGRamseyEightRepOneEfficiency}{32.8\%}
\newcommand{\FGRamseyEightRepTwoEfficiency}{40.6\%}

\newcommand{\FGRamseyOutcome}{All eight cells produced five transactional task winners and a verified final delivery. Median wall time fell from 503.29 seconds at one slot to 153.48 seconds at four slots; the eight-slot median was 153.82 seconds. The frozen graph therefore accelerated up to its ready width of four and showed no additional median improvement when overprovisioned to eight slots. Remote latency variation remained material.}

\newcommand{\FGScriptedOne}{11.41}
\newcommand{\FGScriptedTwo}{6.72}
\newcommand{\FGScriptedFour}{5.91}
\newcommand{\FGScriptedEight}{5.17}
\newcommand{\FGScriptedSixteen}{5.78}

\newcommand{\FGCPStableMedian}{8.984}
\newcommand{\FGCPCriticalMedian}{7.766}
\newcommand{\FGCPOutcome}{All six cells verified. Critical-path ranking selected the long chain first in all three repetitions and reduced median wall time by 13.56\% on this deterministic workload.}

\newcommand{\FGTinyShopTasks}{two in each of two preserved observations}
\newcommand{\FGTinyShopShape}{a serial implementation task followed by a test task (ready width one)}
\newcommand{\FGTinyShopOutcome}{Both executions and final deliveries verified, but both failed the predeclared branch/join golden shape. This is a negative extraction result: correct delivery did not expose root-level speedup.}

\newcommand{\FGJavaScriptTasks}{three}
\newcommand{\FGJavaScriptShape}{two initially ready roots---core implementation and documentation---followed by dependent focused tests}
\newcommand{\FGJavaScriptMakespan}{289.09}
\newcommand{\FGJavaScriptOutcome}{The roots started 9.5 ms apart. Scheduler-owned verification rejected two primary attempts; bounded repairs passed, and final delivery verified all nine Node tests.}


\newcommand{\fgfigure}[3]{%
  \IfFileExists{#1.pdf}{%
    \begin{figure}[t]\centering\includegraphics[width=\linewidth]{#1.pdf}\caption{#2}\label{#3}\end{figure}}{%
    \IfFileExists{#1.png}{%
      \begin{figure}[t]\centering\includegraphics[width=\linewidth]{#1.png}\caption{#2}\label{#3}\end{figure}}{%
      \begin{center}\fbox{\parbox[c][1.0in][c]{0.90\linewidth}{\centering Figure pending\\\texttt{#1}}}\\[-1pt]
      {\footnotesize #2}\end{center}}}}

\newcommand{\fgwidefigure}[3]{%
  \begin{figure*}[t]
    \centering
    \IfFileExists{#1.pdf}{\includegraphics[width=0.88\textwidth]{#1.pdf}}{%
      \IfFileExists{#1.png}{\includegraphics[width=0.88\textwidth]{#1.png}}{%
        \fbox{\parbox[c][1.5in][c]{0.84\textwidth}{\centering Figure pending\\\texttt{#1}}}}}
    \caption{#2}
    \label{#3}
  \end{figure*}}

\title{FirstGreen: HPC-Inspired Parallelism Extraction and Scheduling for Coding Agents}
\author{\FGAuthor}
\date{Technical report for FirstGreen \FGReleaseVersion}

\begin{document}
\maketitle

\begin{abstract}
Coding agents are commonly scaled by opening more independent sessions, but session count is not
application parallelism. Software-engineering goals contain dependencies, write conflicts, serial
verification, and variable-duration work that limit useful concurrency. FirstGreen is a local-first
runtime that applies an HPC perspective to this problem. It extracts a bounded, reviewable task DAG
from a repository and goal; deterministically compiles dependencies, conflicts, resources, and
verification eligibility; estimates work, span, critical path, and exposed width; and schedules
isolated coding-agent attempts under explicit capacity. Verification supplies the completion
boundary required to measure makespan without treating an agent's self-report as success. We
evaluate the same production path with deterministic capacity tests, a controlled-live strong-
scaling matrix, a ready-queue ablation, and Python and JavaScript end-to-end cases. The artifact
publishes frozen inputs, append-only raw outcomes, traces, and negative results. This report asks a
narrow systems question: when a software task exposes safe parallel work, how much of that
parallelism can an agent runtime convert into a shorter time to verified delivery, and where does
the speedup saturate?
\end{abstract}

\section{Introduction}

Running multiple coding-agent sessions in parallel is easy. Deciding which work is independent,
when it is safe to launch, and whether additional capacity can shorten the critical path is not.
A goal may be logically sequential, may contain hidden write conflicts, or may expose several
branches followed by a join. Even a wide graph can fail to scale when repository setup,
verification, integration, or provider latency dominates. Counting agents therefore confuses
allocated capacity with exploitable parallelism.

FirstGreen treats a coding task as a small parallel program. A repository plus an engineering goal
is transformed into semantic work units and then into an approved execution graph. The runtime
uses that graph, rather than independent chat sessions, to make dispatch decisions. Its primary
objective is shorter makespan to a valid delivered result. Reliability mechanisms are necessary
measurement and safety boundaries: without scheduler-owned verification, a fast but incorrect
answer would be recorded as acceleration.

This artifact makes four contributions:

\begin{enumerate}
  \item an end-to-end product path from repository and goal to a bounded, reviewable parallel DAG;
  \item a deterministic compiler that owns graph edges, conflict locks, safety eligibility, and
        work/span analysis even when an LLM proposes semantic work units;
  \item a production scheduler with critical-path ranking, unified capacity and resource admission,
        isolated workspaces, explicit verification, and observable decisions; and
  \item a reproducible evaluation package that separates scripted runtime evidence from live-agent
        scaling and preserves failed, sequential, and negative outcomes.
\end{enumerate}

FirstGreen is not a cluster control plane, an autonomous merge service, or an ML scheduler. The
release targets a local, inspectable systems artifact.

\section{Problem Formulation}

For an approved task DAG $G=(V,E)$, let each task $v$ have a frozen estimated duration $d(v)$.
Total work and critical-path span are

\begin{equation}
  W=\sum_{v\in V}d(v), \qquad
  L=\max_{p\in\mathcal{P}(G)}\sum_{v\in p}d(v).
\end{equation}

$W/L$ is an upper-bound indicator of exposed parallelism, not a speedup promise. For a measured
$N$-slot execution with wall time $T_N$, strong-scaling speedup and efficiency are

\begin{equation}
  S_N=T_1/T_N, \qquad E_N=S_N/N.
\end{equation}

These quantities are meaningful only when each cell replays the same frozen plan, worker model,
verification contract, and stopping condition. FirstGreen defines completion as a verified final
delivery. Failed cells remain failed rather than becoming missing or being replaced after tuning.

The runtime minimizes time to that boundary subject to DAG dependencies, planned write/resource
conflicts, root and subagent capacity, verifier capacity, replay-safety rules, and user approval.
The immediate research questions are:

\begin{itemize}
  \item Can repository-aware extraction find useful branch/join structure while honestly returning
        sequential when it cannot?
  \item Does increasing root capacity reduce makespan for one fixed live DAG, and at what capacity
        does the graph or runtime saturate?
  \item Does bottom-level critical-path ranking improve a scheduler-sensitive workload over stable
        task order without changing the graph?
  \item Does the product path transfer beyond its Python testbed to an independent JavaScript repo?
\end{itemize}

\section{Systems Context}

FirstGreen applies established parallel-systems ideas to agentic software work rather than claiming
a new scheduling theorem. Graham's analysis of precedence-constrained multiprocessor scheduling
shows why adding processors does not imply proportional speedup and can expose timing anomalies
\cite{graham1966}. The work/span vocabulary and emphasis on critical-path depth follow the same
tradition; work-stealing analysis likewise separates total work from the infinite-processor span
\cite{blumofe1999}. FirstGreen differs because tasks are coarse, semantically extracted, modify
shared artifacts, and have stochastic remote execution time. Its deterministic conflict compiler
and verification boundary therefore form part of the schedulable program.

Tail-tolerant services use redundant requests to limit latency variation \cite{dean2013}. FirstGreen
supports bounded delayed hedging for explicitly replay-safe tasks, but disables it in the public
strong-scaling matrix so speculative replicas cannot be confused with parallelism extracted from
the DAG. The present report evaluates a local runtime and does not compare commercial agent
products or claim distributed-cluster scale.

\section{Design}

\fgfigure{figures/system-overview}{FirstGreen's product path. Benchmarking freezes and replays
an approved plan through the same scheduler; it does not introduce a second execution model.}{fig:system}

\subsection{Bounded parallelism extraction}

The scanner produces a read-only, bounded repository map. A semantic planner may propose work
units, artifacts, likely paths, duration hints, and verifier commands, but its proposal is not an
execution plan. Deterministic validation enforces task-count and depth bounds, denied paths,
artifact producers and consumers, acyclicity, and verification eligibility. Overlapping planned
writes compile to capacity-one resources rather than being left as a prompt convention.

The compiled plan records task estimates and their provenance, $W$, $L$, critical paths, initial
ready width, $W/L$, and a hard-bounded root-slot recommendation. The user reviews this material
before approval. A task with insufficient safe parallelism is reported as sequential; the planner
is not forced to invent branches.

\subsection{Scheduling and admission}

The production ready queue supports a stable baseline and bottom-level critical-path rank. Ties
fall back deterministically. Admission considers dependency readiness, root capacity, subagent
budgets, write/resource locks, workspace availability, and shared verifier capacity. Holding a task
on a conflict does not consume a root slot when unrelated ready work can run.

Concurrency policies have hard minima and maxima, a static fallback, and a persisted decision log.
Delayed hedging is separate from failure retry and is disabled unless a task is explicitly replay-
safe. The v0.1 public scaling matrix disables hedging so replica work cannot be mistaken for graph
parallelism.

\subsection{Completion and observability}

Each attempt runs in a dedicated Git worktree; the main working tree is never modified. Scheduler-
owned command and changed-path checks decide whether an attempt is green. A SQLite transaction
selects at most one verified winner. Multi-task sinks are composed into a separate delivery
worktree and verified again. The runtime persists state transitions, ready selections, admission
holds, verifier waits, attempt intervals, sanitized token accounting when available, and final
delivery status. It exports a Chrome/Perfetto trace so idle time and saturation can be explained.

\section{Evaluation Method}

The evidence records source-revision identifiers \texttt{\FGRamseySourceRevision} for Ramsey,
\texttt{\FGJavaScriptSourceRevision} for the JavaScript case, and
\texttt{\FGCPAblationSourceRevision} for the ready-queue ablation. They identify the frozen
pre-release executions but are not reachable commits in the privacy-compacted public history. The
public source is fixed by the immutable release tag \texttt{\FGReleaseRevision}. Experiments were run on \FGHostSummary{} on \FGExperimentDate. Live work
uses \FGPlannerModel{}; reasoning effort is \FGReasoningEffort{}. Exact base commits, issue hashes,
Manifest hashes, command lines, and validity rules are included in the evidence bundle. Its SHA-256
is published in \FGEvidenceBundleChecksum{} to avoid a self-referential archive hash.

The evidence has four deliberately separate layers. First, a deterministic fake-worker branch/join
matrix runs at 1/2/4/8/16 slots to validate scheduler capacity, accounting, worktrees, verification,
delivery, and trace export. Second, the frozen Ramsey workload runs a controlled-live
\FGRamseyRuns{} matrix. Its ready width is four; eight slots test saturation rather than eight-way
work. Third, byte-equivalent stable and critical-path Manifests run three times each with two root
slots. Only ready-queue policy differs. Fourth, TinyShop and an independent Node.js checkout
repository exercise extraction, approval, execution, and delivery as qualitative case studies. The
ablation preserves the frozen one-verifier capacity while replaying two root slots.

The benchmark journal is append-only. $T_1$ is the first valid one-slot observation defined by the
public methodology. No timeout, failed verification, sequential extraction, or interrupted cell is
silently excluded. Provider queue time and financial cost are not claimed when the installed CLI
does not expose them.

\section{Results}

\subsection{Runtime capacity and live strong scaling}

The scripted workload completed in \FGScriptedOne, \FGScriptedTwo, \FGScriptedFour,
\FGScriptedEight, and \FGScriptedSixteen{} seconds at 1, 2, 4, 8, and 16 slots respectively. This
is orchestration evidence, not live coding speed. Its four branches and serial join explain the
efficiency decline after the exposed width is exhausted.

\begin{table}[t]
\centering
\small
\caption{Controlled-live Ramsey raw cells. $T_N$ is wall time in seconds; $S_N$ and $E_N$ use the
first valid one-slot cell as the frozen baseline.}
\label{tab:ramsey}
\begin{tabular}{rrrrr}
\toprule
Slots & Rep. & $T_N$ & $S_N$ & $E_N$ \\
\midrule
1 & 1 & \FGRamseyOneRepOneTime & \FGRamseyOneRepOneSpeedup & \FGRamseyOneRepOneEfficiency \\
1 & 2 & \FGRamseyOneRepTwoTime & \FGRamseyOneRepTwoSpeedup & \FGRamseyOneRepTwoEfficiency \\
2 & 1 & \FGRamseyTwoRepOneTime & \FGRamseyTwoRepOneSpeedup & \FGRamseyTwoRepOneEfficiency \\
2 & 2 & \FGRamseyTwoRepTwoTime & \FGRamseyTwoRepTwoSpeedup & \FGRamseyTwoRepTwoEfficiency \\
4 & 1 & \FGRamseyFourRepOneTime & \FGRamseyFourRepOneSpeedup & \FGRamseyFourRepOneEfficiency \\
4 & 2 & \FGRamseyFourRepTwoTime & \FGRamseyFourRepTwoSpeedup & \FGRamseyFourRepTwoEfficiency \\
8 & 1 & \FGRamseyEightRepOneTime & \FGRamseyEightRepOneSpeedup & \FGRamseyEightRepOneEfficiency \\
8 & 2 & \FGRamseyEightRepTwoTime & \FGRamseyEightRepTwoSpeedup & \FGRamseyEightRepTwoEfficiency \\
\bottomrule
\end{tabular}
\end{table}

\FGRamseyOutcome

\fgwidefigure{figures/ramsey-scaling}{Controlled-live strong scaling for one frozen approved DAG.
Failed or invalid cells are shown rather than imputed.}{fig:scaling}

\fgfigure{figures/ramsey-trace}{Representative scheduler trace showing ready work, root-slot
occupancy, verification, and the serial delivery boundary.}{fig:trace}

\subsection{Scheduling ablation}

The stable policy median was \FGCPStableMedian{} seconds and the critical-path policy median was
\FGCPCriticalMedian{} seconds. \FGCPOutcome{} The causal evidence is dispatch order and the trace,
not only the aggregate time. The policy selected the first long-chain task before both side branches
in all three repetitions and advanced the second chain task before the remaining side branch in two
repetitions.

\fgfigure{figures/critical-path-ablation}{Stable task order versus bottom-level critical-path
ordering on the same scheduler-sensitive DAG.}{fig:ablation}

\subsection{Extraction case studies}

For TinyShop, Luna proposed \FGTinyShopTasks{} tasks with shape \FGTinyShopShape.
\FGTinyShopOutcome{} The case set also includes a sequential goal and a planned write-conflict goal
to demonstrate that extraction quality is not equated with producing the widest possible DAG.

For the independent JavaScript repository, Luna proposed \FGJavaScriptTasks{} tasks with shape
\FGJavaScriptShape. \FGJavaScriptOutcome{} Scheduler makespan was \FGJavaScriptMakespan{} seconds.
This single case supports portability of the product path, not a population-level claim about
JavaScript projects or planner accuracy.

\fgfigure{figures/javascript-live-timeline}{Sanitized independent JavaScript timeline. Failed
primaries do not release dependent work; final delivery follows both repairs.}{fig:javascript-timeline}

\section{Limitations and Threats to Validity}

The controlled-live matrix is small, uses one model, one host, and one fixed workload. Repetitions
capture only limited latency variation. Results do not establish general speedup across repositories,
models, or providers. Task duration estimates influence critical-path ranking but are not trained
predictions. Repository maps are bounded and may omit semantic dependencies. Planned path conflicts
are conservative but cannot prove that two generated patches will integrate cleanly.

Local worktrees isolate edits from the main tree but are not container-grade security boundaries.
Killing a local CLI process may not guarantee server-side cancellation. Verification commands are
repository code and must be treated as potentially dangerous. FirstGreen does not automatically
merge, push, deploy, or execute irreversible external actions.

The extraction case studies are intentionally qualitative. A sequential plan can be a correct
negative result, while a wide plan can be unsafe or too fine-grained. The live scaling result is
conditioned on one already-approved DAG and therefore evaluates scheduling capacity rather than
planner quality. These scopes are separated in every public claim.

\section{Conclusion}

FirstGreen frames coding-agent orchestration as a parallel-systems problem: extract a safe task
graph, quantify work and span, schedule its critical path under real resource constraints, and
measure completion only at verified delivery. Its benchmark harness is a thin observer of that
product path. The public artifact is designed to remain useful if scaling saturates early, a planner
returns sequential work, or a live cell fails: raw outcomes and traces explain where useful
parallelism was lost. The resulting project is an inspectable demonstration of parallelism
extraction, scheduling, isolation, and systems evaluation rather than a multi-session wrapper.

\section*{Artifact Availability}

The tagged source, release assets, evidence archive, reproduction instructions, and checksums are
available at \url{https://github.com/Perilla-desuwa/firstgreen/releases/tag/v0.1.0-rc1}. The release
should be cited using the repository's finalized \texttt{CITATION.cff}.

\bibliographystyle{plain}
\bibliography{references}

\end{document}
